\section{Architettura del Codice}\label{architettura-del-codice}

A livello software, BlueAgent organizza l\textquotesingle interazione
con l\textquotesingle host disaccoppiando l\textquotesingle API in
\passthrough{\lstinline!/api!} dalla gestione a basso livello dei
container e delle componenti di rete, gestita ad alto livello dai
servizi in \passthrough{\lstinline!/services!}, che contiene la logica
di business, mentre le complessità operative sono delegate a moduli e
driver specializzati (\textbf{aggiungi cartelle}).

\begin{center}\rule{0.5\linewidth}{0.5pt}\end{center}

\subsection{\texorpdfstring{Astrazione del Runtime:
\texttt{ContainerEngine}}{Astrazione del Runtime: ContainerEngine}}\label{astrazione-del-runtime-containerengine}

BlueAgent non invoca direttamente il demone Incus
all\textquotesingle interno della logica delle challenge. Tutte le
operazioni di gestione del ciclo di vita dei container sono intermediate
da un\textquotesingle interfaccia astratta denominata
\passthrough{\lstinline!ContainerEngine!}.

Questa classe astratta formalizza esclusivamente le primitive necessarie
ai requisiti operativi del Cyber Range:

\begin{itemize}
\tightlist
\item
  Creazione, avvio, arresto forzato ed eliminazione delle istanze.
\item
  Duplicazione rapida e gestione delle interfacce di rete.
\item
  Esecuzione dei comandi nei container.
\end{itemize}

Disaccoppiare l\textquotesingle interfaccia
dall\textquotesingle implementazione permette ai servizi di non dover
conoscere come avviene la gestione dei container, e di predisporre il
codice per poter utilizzare altri gestori.

\begin{quote}
\textbf{ATTENZIONE: anche la generazione delle immagini dovrebbe
avvenire tramite interfaccia}
\end{quote}

L\textquotesingle implementazione concreta
\passthrough{\lstinline!IncusEngine!} comunica con il demone di sistema
tramite Unix Domain Socket
(\passthrough{\lstinline!/var/lib/incus/unix.socket!}), traducendo le
primitive astratte in chiamate REST HTTP asincrone.

\begin{quote}
\passthrough{\lstinline!Unix Domain Socket (UDS)!}: Meccanismo di
Inter-Process Communication (IPC) gestito interamente dal kernel del
sistema operativo. La comunicazione avviene trasferendo i dati
direttamente tra i buffer di memoria dei processi nello spazio kernel
senza attraversare lo stack di rete. L\textquotesingle endpoint è
rappresentato da un file speciale sul filesystem (es:
\passthrough{\lstinline!/var/lib/incus/unix.socket!}).\footnote{\label{fn:ref_21}Linux
  manual page, unix(7) - Sockets for local interprocess communication:
  \url{https://man7.org/linux/man-pages/man7/unix.7.html}}
\end{quote}

\begin{center}\rule{0.5\linewidth}{0.5pt}\end{center}

\subsection{\texorpdfstring{Ciclo di Vita della Sandbox: Async Context
Manager
(\texttt{OVNVerifier})}{Ciclo di Vita della Sandbox: Async Context Manager (OVNVerifier)}}\label{ciclo-di-vita-della-sandbox-async-context-manager-ovnverifier}

Il processo di verifica richiede di istanziare una replica fedele della
rete di produzione che deve esistere esclusivamente per la durata dei
test. In caso di qualsiasi tipo di fallimento,
l\textquotesingle infrastruttura clonata deve essere rilasciata per
evitare risorse allocate e sprecate.

Per garantire questa proprietà, il ciclo di vita della sandbox è
incapsulato all\textquotesingle interno della classe
\passthrough{\lstinline!OVNVerifier!}, implementata come
\textbf{Asynchronous Context Manager} di Python.

\begin{quote}
\textbf{Pattern RAII (Resource Acquisition Is Initialization):}
L\textquotesingle allocazione delle risorse è vincolato
dall\textquotesingle acquisizione del contesto, mentre il loro rilascio
e smantellamento deterministico è garantito all\textquotesingle uscita
dal blocco, gestendo automaticamente qualsiasi anomalia o interruzione.
\footnote{\label{fn:ref_22}Bjarne Stroustrup, Resource Acquisition Is Initialization
  (RAII), C++ Programming Language.}
\end{quote}

Il ciclo di vita dell\textquotesingle{}
\passthrough{\lstinline!OVNVerifier!} segue tre fasi:

\begin{enumerate}
\def\labelenumi{\arabic{enumi}.}
\tightlist
\item
  \textbf{Inizializzazione (\passthrough{\lstinline!\_\_aenter\_\_!}):}

  \begin{itemize}
  \tightlist
  \item
    Creazione della rete OVN temporanea
  \item
    Duplicazione dei container tramite snapshot ZFS
  \item
    Riconnessione dell\textquotesingle interfaccia
    \passthrough{\lstinline!eth0!} dei cloni alla nuova rete OVN
  \item
    Istanziazione opzionale del container \emph{Verifier Runner}.
  \end{itemize}
\item
  \textbf{Esecuzione delle Verifiche
  (\passthrough{\lstinline!run\_checks!}):}

  \begin{itemize}
  \tightlist
  \item
    Esecuzione controllata dei comandi di test con timeout
  \end{itemize}
\item
  \textbf{Rilascio e Teardown
  (\passthrough{\lstinline!\_\_aexit\_\_!}):}

  \begin{itemize}
  \tightlist
  \item
    Smantellamento delle risorse:

    \begin{enumerate}
    \def\labelenumii{\arabic{enumii}.}
    \tightlist
    \item
      Arresto ed eliminazione del container
      \passthrough{\lstinline!verify-runner!};
    \item
      Arresto forzato ed eliminazione di tutti i container cloni;
    \item
      Eliminazione della rete virtuale OVN.
    \end{enumerate}
  \end{itemize}
\end{enumerate}

L\textquotesingle uso dell\textquotesingle{}\textbf{Async Context
Manager} fornisce dunque un metodo sicuro per utilizzare
l\textquotesingle oggetto \passthrough{\lstinline!OVNVerifier!}

\begin{center}\rule{0.5\linewidth}{0.5pt}\end{center}

\subsection{Concorrenza e Two-Tier
Locking}\label{concorrenza-e-two-tier-locking}

Incus non supporta operazioni concorrenti sul medesimo container.
Inoltre per evitare problemi di accesso alle risorse e\textquotesingle{}
opportuno implementare un sistema di lock.

BlueAgent implementa un meccanismo di sincronizzazione a due livelli
(\passthrough{\lstinline!TwoTierLock!}), dove sequenzialmente si
richiede un tier dopo l\textquotesingle altro:

\begin{itemize}
\tightlist
\item
  \textbf{Tier 1 (Intra-processo):} utilizza
  \passthrough{\lstinline!asyncio.Lock!} per effettuare il lock delle
  risorse all\textquotesingle interno di un singolo processo.
\item
  \textbf{Tier 2 (Inter-processo):} applica un lock non bloccante sul
  filesystem (\passthrough{\lstinline!fcntl.flock!}) condiviso tra piu
  processi
\end{itemize}

Questo lock protegge l\textquotesingle accesso alle risorse in due casi:

\begin{itemize}
\tightlist
\item
  BlueAgent esegue su un solo processo (worker), ma le richieste
  potrebbero essere gestite da piu task contemporaneamente. Allora il
  lock Tier 1 protegge l\textquotesingle accesso alla stessa risorsa da
  parte di due task sullo stesso worker.
\item
  BlueAgent esegue su piu processi (worker), allora il lock Tier 2
  protegge l\textquotesingle accesso alla stessa risorsa da parte di due
  task su worker differenti.
\end{itemize}

\begin{quote}
Come si avviano piu\textquotesingle{} worker?
\passthrough{\lstinline!uvicorn main:app --workers 4!}
\end{quote}

\begin{quote}
\passthrough{\lstinline!uvicorn!}: e\textquotesingle il web server su
cui si basa FastAPI, dunque e\textquotesingle{} uvicorn che decide come
gestire le richieste, e a quali task e worker assegnarli. Da solo
FastAPI non funziona, non implementa il server di rete.
\end{quote}

Questa architettura protegge le risorse sia a livello di singolo
processo sia in caso di server multi-worker, garantendo mutua esclusione
sulle mutazioni su tutti i tipi di risorse.

La risorsa principale su cui avviene la mutua esclusione nel codice
e\textquotesingle{} il team. Qualsiasi operazione sul team \(k\) genera
un lock a 2 tier per quel team.

\begin{center}\rule{0.5\linewidth}{0.5pt}\end{center}

\subsection{Pipeline Asincrona non Bloccante su Socket
Unix}\label{pipeline-asincrona-non-bloccante-su-socket-unix}

Le operazioni di su container possono comportare latenze
nell\textquotesingle ordine di decine o centinaia di millisecondi.
L\textquotesingle adozione di un modello sincrono sarebbe totalmente
inefficiente.

BlueAgent è progettato secondo un paradigma completamente asincrono in
modo da adattarsi al metodo \passthrough{\lstinline!wait!} di Incus:

\begin{itemize}
\tightlist
\item
  \passthrough{\lstinline!IncusEngine!} comunica con il socket UDS
  locale (\passthrough{\lstinline!/var/lib/incus/unix.socket!}) per
  mezzo di \passthrough{\lstinline!httpx.AsyncClient!}.
\item
  Poiché le operazioni bloccanti di Incus sono gestite in background,
  BlueAgent interroga asincronamente l\textquotesingle endpoint
  \passthrough{\lstinline!/1.0/operations/<uuid>/wait!} per sapere
  quando l\textquotesingle operazione viene rilasciata.
\end{itemize}

Inoltre l\textquotesingle esecuzione dei comandi sulla macchina host
avviene tramite:
\passthrough{\lstinline!asyncio.create\_subprocess\_exec!} per evitare
chiamate bloccanti su comandi IO intensive. (CONTROLLA SE HO DETTO UNA
MINCHIATA)

Questo approccio consente a un singolo processo BlueAgent di servire le
richieste di telemetria utilizzando piu task e senza rimanere in attessa
su una singola richiesta evitando di servire le altre richieste.

\begin{center}\rule{0.5\linewidth}{0.5pt}\end{center}

\subsection{\texorpdfstring{Resilienza e Riconciliazione dello Stato
(\texttt{LifecycleManager})}{Resilienza e Riconciliazione dello Stato (LifecycleManager)}}\label{resilienza-e-riconciliazione-dello-stato-lifecyclemanager}

Qualora l\textquotesingle host subisca un riavvio forzato improvviso o
il processo di BlueAgent venga terminato bruscamente con
\passthrough{\lstinline!SIGKILL!}, le routine di pulizia
(\passthrough{\lstinline!\_\_aexit\_\_!}) non possono intervenire,
lasciando container e switch virtuali orfani sul sistema.

Per risolvere queste anomalie senza intervento manuale, BlueAgent
include il demone in background
\passthrough{\lstinline!LifecycleManager!}:

\begin{itemize}
\tightlist
\item
  Scansiona periodicamente le istanze e le reti con prefisso
  \passthrough{\lstinline!verify-*!} presenti su Incus, forzandone
  l\textquotesingle eliminazione se risultano inattive.
\item
  \textbf{Grace Period di 90 Secondi:} il demone ignora qualsiasi
  risorsa la cui creazione risalga a meno di 90 secondi. Questa finestra
  temporale risolve una race condition fondamentale: durante
  l\textquotesingle avvio di una verifica, la rete OVN viene creata
  prima che i container cloni vi siano collegati.
\end{itemize}

\begin{center}\rule{0.5\linewidth}{0.5pt}\end{center}
